% This must be in the first 5 lines to tell arXiv to use pdfLaTeX, which is strongly recommended.
\pdfoutput=1
% In particular, the hyperref package requires pdfLaTeX in order to break URLs across lines.

\documentclass[11pt]{article}

% Remove the "review" option to generate the final version.
\usepackage[review]{acl}

% Standard package includes
\usepackage{times}
\usepackage{latexsym}

% For proper rendering and hyphenation of words containing Latin characters (including in bib files)
\usepackage[T1]{fontenc}
% For Vietnamese characters
% \usepackage[T5]{fontenc}
% See https://www.latex-project.org/help/documentation/encguide.pdf for other character sets

% This assumes your files are encoded as UTF8
\usepackage[utf8]{inputenc}

% This is not strictly necessary, and may be commented out,
% but it will improve the layout of the manuscript,
% and will typically save some space.
\usepackage{microtype}

% If the title and author information does not fit in the area allocated, uncomment the following
%
%\setlength\titlebox{<dim>}
%
% and set <dim> to something 5cm or larger.

\title{Can Preprocessing and Hyperparameter Tuning Close the Gap? A Study of Different Models for Tweet Sentiment Classification}

% Author information can be set in various styles:
% For several authors from the same institution:
% \author{Author 1 \and ... \and Author n \\
%         Address line \\ ... \\ Address line}
% if the names do not fit well on one line use
%         Author 1 \\ {\bf Author 2} \\ ... \\ {\bf Author n} \\
% For authors from different institutions:
% \author{Author 1 \\ Address line \\  ... \\ Address line
%         \And  ... \And
%         Author n \\ Address line \\ ... \\ Address line}
% To start a seperate ``row'' of authors use \AND, as in
% \author{Author 1 \\ Address line \\  ... \\ Address line
%         \AND
%         Author 2 \\ Address line \\ ... \\ Address line \And
%         Author 3 \\ Address line \\ ... \\ Address line}

\author{A0251997L, A0252459X, A0257425B, A0252829W, A0252183J \\
  Group 28 \\
  Mentored by Esther Gan \\
  \texttt{\{e0957585,e0958047,e0968397,e0958417,e0957771\}@u.nus.edu}}

\begin{document}
\maketitle
\begin{abstract}
Sentiment classification in social media is complicated by informal language, negation, and contextual ambiguity. 
While modern Transformer-based models define the state-of-the-art (SOTA), classical algorithms like Naive Bayes (NB) remain relevant due to their computational efficiency and interpretability. 
This paper investigates the extent to which meticulous preprocessing, feature engineering, and hyperparameter tuning can compensate for the structural limitations of the Bag-of-Words (BoW) assumption. 
We implement an empirical progression starting with a highly optimized Multinomial Naive Bayes pipeline—incorporating negation tagging, emoticon translation, and VADER lexicon integration—achieving a competitive accuracy of 71\% and a macro-F1 of 0.70.

We then evaluate these same features within a Bidirectional LSTM (Bi-LSTM) and a fine-tuned RoBERTa model to identify at what point manual feature engineering becomes redundant or counterproductive due to superior representation learning. 
A qualitative error analysis reveals that while NB performs robustly on high-density sentiment tokens, it consistently fails in cases of sentiment inversion, polysemy, and spurious neutral correlations—areas where RoBERTa`s self-attention mechanism excels. 
Our results suggest that while careful normalization significantly closes the performance gap in simpler datasets, architectural context-awareness remains indispensable for capturing the special nuances found in tweets.
\end{abstract}

\section{Introduction}

Sentiment analysis is a fundamental task in natural language processing (NLP), with applications ranging from market analysis to public opinion monitoring. Social media platforms such as Twitter present a particularly interesting setting for sentiment classification due to the informal, noisy, and highly contextual nature of the language used. Tweets often contain misspellings, abbreviations, emojis, and linguistic constructs such as negation that can significantly affect sentiment interpretation. These characteristics make Twitter sentiment classification a challenging benchmark for machine learning models.

Traditional machine learning approaches such as Naive Bayes have long been used for text classification tasks due to their simplicity, efficiency, and relatively strong performance on bag-of-words representations. However, Naive Bayes relies on strong independence assumptions such as the Bag of Words (BoW) assumption, which limits its ability to capture linguistic phenomena such as negation, contextual dependencies, and word interactions. 
Sequential neural models such as Bidirectional LSTMs address these limitations by learning order-sensitive contextual representations, but remain constrained by the amount of in-domain training data available.
In contrast, modern pretrained transformer models such as RoBERTa have become state-of-the-art(SOTA) for sentiment analysis by leveraging large-scale pretraining to obtain rich contextual representations that generalise well to downstream tasks with minimal fine-tuning.

Despite the superior performance of transformer-based models, classical algorithms like Naive Bayes remain attractive due to their interpretability, computational efficiency, and ease of deployment. 
A key question is therefore how far careful preprocessing and feature engineering compensate the structural limitation of simpler models.
For Naive Bayes, preprocessing techniques and feature engineering compensates for model capability and strong BoW assumptions, for example explicit negation tagging, handling of expressive lengthening and lexicon injections.

Through this comparison, we aim to identify the extent to which preprocessing can compensate for the structural limitations of Naive Bayes, and to highlight the types of linguistic patterns that remain difficult for such models to capture. 
Then we evaluate if the same preprocessing and engineered feature become redundant or countereffective for Bi-LSTM as the improved architecture learns it equivalent from sequential context and better token representation with GloVe. 
Lastly, RoBERTa, where the preprocessing and feature engineering might be disruptive due to its superior architecture and large scale pretraining.

We investigate this progression empirically. We first build a Naive Bayes pipeline with extensive preprocessing and feature engineering to push the limits of the Naive Bayes. We then apply the same preprocessing to a Bi-LSTM to evaluate whether steps that helped 
Naive Bayes remain beneficial under a sequential architecture, supported with better word representation (GloVe). Then finally, on RoBERTa, to see if the remaining feature engineering and processing was beneficial to it with large scale pre training. Through this, we can also uncover remaining weaknesses in Naive Bayes, by 
training Naive Bayes and RoBERTa on the same dataset to inspect instances where RoBERTA succeeded while Naive Bayes failed.
%%%%%

\section{Related Work / Background}

\subsection{State of the Art Models}

There has been much research in the field of sentiment analysis for web content such as tweets. Most of the top performing models use LSTM-type models\citep{ruder-2026}. 

According to \url{https://nlpprogress.com}, the best sentiment analysis model (up to 2019)\citep{yang-2019} had an accuracy of 96\%, using a self-attention setup. The only preprocessing done was just encoding relative positioning in the vector representation.

However, as transformers are more computationally expensive and harder to interpret, we wanted to see if the standard Naive Bayes with preprocessing was enough to rival the performance of the best models.

\subsection{Evaluating the Effectiveness of Text Pre-Processing in Sentiment Analysis (2022) by Palonimo et al.}

However, significant research has also been done in the industry on improving preprocessing methods.

In this paper\citep{palomino-2022}, where they showcased a Naive Bayes Classifer used on tweets, and tested various preprocessing pipelines to test the effectiveness of diifferent preprocessing pipelines. They have also provided a table of their accuracy scores they obtained with different preprocessing pipelines, as well as with no preprocessing, which allows for a good benchmark for our models.

In both cases, it was shown that Naive Bayes managed to obtain a very high accuracy, with as high as 91\% on gold standard datasets, and 86\% on the whole dataset overall. This highlights its potential to be a decent classifier despite its simplicity. 91\% is also very close to the self-attention model's accuracy, at 96\%, suggesting that it might be able to come close to the accuracy of more sophisticated models.

However, the tweets they used were from the political domain, about the Grand Old Party (GOP) held among the candidates of the Republican Party for the 2016 US Presidential Election. Naturally, this makes the tweets much more polarizing, making it easier for those models to achieve better scores.

With our much more neutral and nuanced dataset from Kaggle, our goal is to see if Naive Bayes could still hold up against a dataset with not so polarizing comments.

\subsection{Implementation of Naive Bayes Algorithm on Sentiment Analysis Application (2021) by Martiti et al.}

In this paper\citep{martiti-2021}, there is yet another sentiment analysis classifier made with simple Naive Bayes, but this time done on Grab Twitter comments. This is done with TF-IDF vectors instead of Count vectors. As shown, they managed to achieve an impressive 99.623\% accuracy on their Naive Bayes classifier, which highlights the massive potential in just using Naive Bayes.

However, as the dataset is different (Tweet comments about Grab instead of general tweets), which are usually more polarizing by nature, it might not reflect the actual potential of just using Naive Bayes on more neutral tweets. Furthermore, as our tweets cover a much wider range of context, rather than just one (Grab), it will likely be more difficult to just use Naive Bayes on the tweet dataset. As such, we want to try to push Naive Bayes as far as we can with our dataset.

\subsection{Summary Comments on Industry Projects}

From these industry projects, we thus have hope that we can simply push Naive Bayes to its limit in doing sentiment analysis on tweets, and thus will be exploring ways to improve the accuracy of Naive Bayes with preprocessing methods.

\section{Corpus Analysis \& Method}

The dataset consists of tweet texts labelled with one of three sentiment classes: negative, neutral, or positive. 
The raw training set contains 27,481 rows and the raw test set contains 4,815 rows. 
After removing rows with missing values in critical fields such as text and sentiment, 1 row was dropped from the training set and 1,281 rows were dropped from the test set, leaving 27,480 training rows and 3,534 test rows. 
The distribution of sentiment in the training dataset is as follows.
\begin{table}[!ht]
\centering
\begin{tabular}{lcr}
\hline
\textbf{Sentiment Label} & \textbf{Count} \\
\hline
Negative & 7,781 \\
Neutral & 11,117 \\
Positive & 8,582 \\
\hline
\end{tabular}
\caption{Distribution of sentiment labels in the training dataset}
\end{table}
As the distribution of sentiment seems to be rather balanced with no significant minority, we chose not to apply class-balancing techniques as oversampling and undersampling incur their own tradeoffs.

In addition to tweet text, the dataset includes metadata such as country, age group, time of tweet, population, land area, and density. 
However, our exploratory analysis suggested that these metadata features highlighted that sentiment distribution were largely balanced across these categories, hence are weak predictors of sentiment. 
Therefore, we chose to drop these auxiliary variables and focus primarily on the tweet text itself. Furthermore, the inclusion of these variables would have made the interpretability of the models more difficult, as it would be less clear how these features interact with the text features to influence sentiment predictions.

A `selected\_text' column, which is a subset of the value present in `text' was provided in the train datatset and not the test dataset.
While the process of deriving and purpose of `selected\_text' is not explicitly documented in the dataset, we inferred that it represents the most relevant part of the tweet that led to the assignment of the sentiment label of the text. 
Since this column was not provided in test set, we are unable to validate it and hence excluded from training and testing purposes.

As our goal is to identify preprocessing steps that bridges the gap between Naive Bayes, we consider and begin with preprocessing and feature engineering that would most improve the performance of Naive Bayes.

\subsection{Preprocessing}
\begin{enumerate}
    \item \textbf{Removal of rows with critical missing values.} Rows with missing values in the ``text'' and ``sentiment'' columns are removed. It was manually observed that all the testing and training rows that are dropped are merely empty rows.
    \item \textbf{Lowercasing.} Converting tokens to lowercase reduces vocabulary sparsity. Although capitalization, especially all caps (e.g., ``I am SO happy''), can carry sentiment information, we found that the majority of capitalized tokens in the dataset were either at the beginning of sentences or were proper nouns, which are not strong indicators of sentiment. Therefore, we opted to lowercase all tokens to improve generalization and reduce noise.
    \item \textbf{Removal of URLs.} URLs are removed as they are not informative for sentiment classification and introduces noise into the feature space.
    \item \textbf{Removal of numbers.} Numeric tokens are removed because they rarely carry stable sentiment signals in the corpus and can introduce noise into bag-of-words features (e.g., for the sentence ``I earned 100 dollars after a day's work'', it is uncertain whether 100 is a large enough number for this sentence to be positive or negative). 
    
    \item \textbf{Lemmatisation.} Lemmatisation reduces the vocabulary size and therefore improve the frequency estimation for rare forms. However, the morphological form may carry sentiment relevant signals and removing them might affect performance, especially if the dataset is big enough to give sufficiently high frequencies for rare forms. We evaluated the performance with and without lemmatisation, eventually deciding to keep the morphological forms as lemmatisation in fact resulted in accuracy drops. 
    \item \textbf{Stopword Removal.} Removing Stopwords reduces the vocabulary size but is widely known to risk performance as by removing critical words that are negative in nature (e.g `not', `should'nt'). Therefore, we curated a list of certain stopwords that are crucial and are not excluded during the removal process (e.g not, shouldn't etc). 
\end{enumerate}

\subsection{Feature Engineering}
\begin{enumerate}
    \item \textbf{Translation of emoticons.} Emoticons are translated to their text equivalent ((e.g :( -> `sad face' and :) -> `happy face')) to prevent them from getting regarded as OOV (especially since they may be fragmented by tokenisation). This provides them more sentiment signal for improved performance.
    \item \textbf{Removal of expressive lengthening.} Expressive lengthening such as ``soooo'' could be replaced with ``soo'' to reduce the vocabulary space. The logic involves replacing repeated characters of more than length 2 with only 2 characters. Our experiments indicate that performance gains caused by a smaller vocabulary space did not outweigh the loss of sentiment information caused by the removal of expressive lengthening, indicating that the expressive lengthenings are meaningful indicators of sentiment.
    \item \textbf{Injection of VADER Sentiment Token.} Valence Aware Dictionary and Entiment Reasoner (VADER) \citep{hutto-2014} is a lexicon and rule based sentiment tool trained on social media text. Given a text, outputs a score from -1 to 1, to which we used to append a sentence level sentiment token `VADER\_POSITIVE', `VADER\_NEGATIVE' and `VADER\_NEUTRAL' at the end of the text if the score is \(\geq 0.05 \), \(\leq 0.05 \) or otherwise respectively. 
    This partially proxies the large training corpus that RoBERTa and LSTM are trained on, by granting the model gains access to an external sentiment signal derived from a rich, human-validated lexicon.
    \item \textbf{Injection of Fragmented VADER Sentiment Token.} While similar to the above, there is a first split of the text on clause boundaries (e.g `,', `;', `but') then apply VADER on each fragments to append tokens such as `FRAG1\_NEGATIVE' and `FRAG2\_NEUTRAL' to implied the sentiment of each segment. This better represenation of polarity shifts that Naive Bayes may miss as a result of the independence assumption it holds, that LSTM and RoBERTa does not face due to it ability to handle long ranges sequential data. 
    \item \textbf{Negation Tagging.} How it helps bridge the gap: Naïve Bayes treats each token as conditionally independent, so ``not good" is represented as two unrelated features, not effectively capturing the polarity shift down the sentence. This method addresses that by scanning tokens sequentially: when a negation word (e.g., not, never, no) is encountered, all subsequent tokens are prefixed with `NOT\_' until a punctuation mark or contrast word (e.g `but', `however', `although') or when the `negation\_window' (if not 0) is reached. 
    Additionally, a `CONTRAST\_SHIFT' token is added when a contrast word is added. The `NOT\_' prefix creates a distinct vocabulary that encodes the negatation directly into the token, while the contrast shift token serves as a signal to the model of a polarity shift in the text.
    \item \textbf{Injection Subject-Verb-Object.} One of the core advantages of RoBERTa and LSTMs is their ability to model relationships between non-adjacent tokens, which Naive Bayes does not have due to its independent assumption. Therefore, we attempt to proxy it by identifying and appending subject-verb-object tokens. spaCy's dependency parser is used to extract the ROOT verb along with its nominal subjects and objects/attributes, then a compact token such as `SVO\_i\_love\_movie' is appended to the text.
    \item \textbf{Polarity Propagation.} Similar to relying on a quality and large external corpus, we manually crafted a sequential scoring pass over the preprocessed text (after negation tagging). It accumulates a signed polarity score by awarding \(\pm 2\) for positive/negative lexicon matches, flipping the sign on `NOT\_' prefixed tokens, halving the score at a `CONTRAST\_SHIFT' boundary, and applying a 1.5x boost if any all-caps token was detected (a common intensity marker in informal text). The final score is discretized into `POLARITY\_STRONG\_POS', `POLARITY\_STRONG\_NEG', or `POLARITY\_NEUTRAL' and appended as a token. This also proxies the long-range dependencies issues Naive Bayes faces.
    \item \textbf{Normalising Expressive Lengthening.} The normalisation of expressive lengthening (eg. gooood) is crucial as they carry strong sentiment signal, but are often discarded due to it being treated differently from its root word (good). Therefore, normalisation is essential to identify and preserve the token, moreover, it reduces the vocabulary size that will be unnecessarily large without normalisation.
\end{enumerate}

\section{Experiments}
Naive Bayes will serve as our baseline model with all the preprocessing steps to push the limits of performance and interpretability. Ablations of the preprocessing steps and feature engineered is done after to identify methods that may not
be as useful for Naive Bayes and plausible rationales. Then Bi-LSTM will be applied with the preprocessing steps and feature engineering that served the Naive Bayes model well to evaluate whether the contribution to the model is counterproductive. 
For feature engineered and preprocessing steps that were counterproductive, we will attempt to identify the rationale and highlight the relationship between explicit preprocessing and feature engineering with better model architecture and better representation learning.
We finally contrast it with RoBERTa, to showcase what the upper limit of the SOTA can achieve as well as whether certain preprocessing and engineered feature that may have served the Bi-LSTM well does not contribute as much to RoBERTa along with the ratonale.
The code and result can be found on our GitHub repository\footnote{\url{https://github.com/jinyang628/cs4248}}.


\subsection{Naive Bayes method}
We trained a baseline Multinomial Naive Bayes classifier with basic preprocessing such as lowercasing, removal of URLs, removal of numbers and translating emoticons. We performed a hyperparameter grid search over the vectorizer (TF-IDF or CBoW) and model settings, including the n-gram range, minimum document frequency, vocabulary size, and smoothing parameter alpha, to identify the best
set of hyperparamters that we will used across the ablations test of incrementally adding the carefully designed preprocessing and engineered features, to ensure complexity and reduce complexity.

The best-performing configuration and the test ranges used are presented in Table 1.

\begin{table}[!ht]
\centering
\begin{tabular}{lcr}
\hline
\textbf{Hyperparameter} & \textbf{Value} & \textbf{Test Range} \\
\hline
Max Features & 4000 & 1000-30000 \\
N-gram Range & (1,2) & (x,y); $x,y \in $ [1,5] \\
Minimum df & 4 & 1-5 \\
Alpha & 1.0 & 0.1-1.0\\
Vectorizer & Count & Count, TF-IDF \\
\hline
\end{tabular}
\caption{Best-performing Naive Bayes hyperparameters}
\end{table}

The optimal n-gram range indicates that unigram and bigram features were sufficient, while more complex n-gram ranges did not yield better cross-validation performance in the saved notebook.

Unlike standard multiclass classification, our notebook maps sentiment labels to ordinal scores {-1, 0, 1} and also experiments with threshold calibration for converting prediction scores back to discrete classes. This was done to improve macro-F1, especially for the neutral class, which often lies between clearly positive and clearly negative examples.

\subsection{Naive Bayes results}

\begin{table}[h]
\small
\centering
\renewcommand{\arraystretch}{1.2}
\begin{tabular}{p{5.8cm} c c}
\hline
\textbf{Preprocessing} & \textbf{Accuracy} & \textbf{Macro F1} \\
\hline
Baseline & 0.66 & 0.66 \\
\hline
+ Translate Emoticon               & 0.66 & 0.66 \\
+ Normalise Elongation             & 0.66 & 0.66 \\
+ SVO Triplets                     & 0.66 & 0.66 \\
+ Position Tagging                 & 0.66 & 0.66 \\
\hline
+ VADER + Frag VADER               & 0.70 & 0.70 \\
+ Negation Scoping                 & 0.67 & 0.66 \\
+ Polarity Propagation             & 0.67 & 0.67 \\
+ Lemmatise                        & 0.67 & 0.66 \\
\hline
\textbf{+ VADER + Frag VADER}      &      & \\
\textbf{+ Negation Scoping}        &      & \\
\textbf{+ Polarity Propagation + Lemmatise} & \textbf{0.71} & \textbf{0.71} \\
\hline
\end{tabular}
\caption{Naive Bayes Ablation Study}
\label{tab:nb_ablation}
\end{table}

The baseline Naive Bayes model, with the best set of hyperparameters, obtained an accuracy of 0.66 and Macro F1 score of 0.66, which is a decent classification model performing better than random guessing.
The engineered features and preprocessing such as `Translate Emoticon` and Normalise Elongation did not improve the performance of the model.
In contrast, the five engineered features and preprocessing, collectively contributed substantially to performance of the model. 
This is suggest that VADER served as a meaningful substitute for the pre-trained semantic grounding that RoBERTa and LSTM (GloVe) has that NB cannot derive from co-occurrence counts alone. This produced the single largest individual gain of +0.04, confirming that the absence of external lexical knowledge is one of the most significant bottlenecks for NB on sentiment tasks.
Negation scoping also contributed, serving as effective signals for the propagation of negation mitigating the independence assumption of Naive Bayes.
Similarly, Polarity propagation, effectively approximate long range dependencies of text, accumulating the sentiment across the text and encoding it as a single token that served as an effective sentiment signal.
Lastly, lemmatisation expectedly contributed to the performance by reducing inflected variants into the same word, reducing vocabulary sparsity and improving the reliability of the frequency estimates that NB depends on.

This suggests that by designing features and preprocessing around the limitations of Naive Bayes and strength of complex models like Bi-LSTM and RoBERTa, we can significantly improve the performance of a simple model.

\subsection{LSTM method}

We implemented a Bidirectional Long Short-Term Memory (Bi-LSTM) classifier as an intermediate model that bridges the link from Naive Bayes to RoBERTa. 
The choice of choice of Bi-LSTM is motivated by two key limitations of Naive Bayes, 
the bag-of-words independence assumption renders it unable to model word order or sequential context and it cannot capture long-range dependencies, which are critical for sentiment in compound and complex sentences.
A Bi-LSTM processes the token sequence both left-to-right and right-to-left, learning contextual representations that encode the influence of surrounding tokens, and is therefore structurally better 
suited to these phenomena. The Bi-LSTM is trained with GloVe a pre trained embedding to 

We trained a baseline Bi-LSTM with basic preprocessing such as lowercasing, removal of URLs, removal of numbers and translating emoticons. 
We performed a hyperparameter grid search over model parameters such as number of hidden dimension and number of layers to identify the best
set of hyperparamters that we will used across the ablations test of incrementally testing if the preprocessing and engineered features that helped Naive Bayes remains useful for the Bi-LSTM.

The best-performing configuration and the test ranges used are presented below.
\begin{table}[h]
\centering
\begin{tabular}{lc}
\hline
\textbf{Hyperparameter} & \textbf{Value} \\
\hline
Hidden dimensions & [256] \\
Dropout & 0.2  \\
Batch size & 64 \\
Number of layers & 1 \\
Bidirectional & True \\
\hline
\end{tabular}
\caption{Best-performing LSTM hyperparameters (Embedding dimension fixed at 100 to align with embedding dim of GloVe)}
\end{table}

\subsection{Bi-LSTM results}

\begin{table}[h]
\centering
\small
\renewcommand{\arraystretch}{1.2}
\begin{tabular}{p{4.5cm} c c}
\hline
\textbf{Preprocessing} & \textbf{Accuracy} & \textbf{Macro F1} \\
\hline
Baseline & 0.72 & 0.72 \\
\hline
+ Lemmatise                 & 0.75 & 0.75 \\
\hline
+ Lemmatise + VADER         &      & \\
\quad + Frag VADER + Negation Scoping        &      & \\
\quad + Polarity Propagation                 & 0.75 & 0.75 \\
\hline
\end{tabular}
\caption{Bi-LSTM Ablation Study}
\label{tab:lstm_ablation} 
\end{table}

The Baseline Bi-LSTM without any of the engineered features and preprocessing that were beneficial to Naive Bayes already outperformed the best Naive Bayes.
This showcase the superiority of a more complexed architecture and strong pre trained embeddings. Except for lemmatisation, the addition of the engineered features and preprocessing
beneficial to Naive Bayes did not contribute to an improve in performance. This is largely expected since LSTM's sequential model could better represent negation and polarity shift than the engineered features.
Lemmatisation significantly improved the model as without it, the vocabulary coverage by the GloVe embedding is around 40\%, but its introduction significantly increased coverage to around 73\%, further illustrating the value of a large scale pre training corpus on model performance.
Overall, this suggests that the efforts of a carefully designed feature engineering and pre processing step, may generally be outweighted by a better model architecture and larger and higher quality pre train corpus.

\section{RoBERTa method}

We fine-tuned DistilRoBERTa to highlight the importance of a better model architecture and pre training corpus on its ability to better represent and identify sentiment, over the efforts to improve on the Bi-LSTM and Naive Bayes. 
Instead of directly treating the task as a 3-way classification problem, the notebook formulates it as regression over sentiment scores, where negative, neutral, and positive are mapped to -1, 0, and 1 respectively. 
The regression outputs are then converted back into class labels using thresholding.
Basic preprocessing such as lowercasing, removal of URLs, removal of numbers and translating emoticons, were applied.
The model used a maximum sequence length of 64, batch size of 32, learning rate of 2e-5, weight decay of 0.01, and one training epoch. This design keeps training lightweight while still leveraging contextual representations learned from large-scale pretraining.
\subsection{RoBERTa results}

\begin{table}[h]
\centering
\caption{RoBERTa Ablation Study}
\renewcommand{\arraystretch}{1.2}
\begin{tabular}{p{3cm} c c} 
\hline
\textbf{Preprocessing} & \textbf{Accuracy} & \textbf{Macro F1} \\
\hline
RoBERTa  & 0.78 & 0.78 \\
RoBERTa + Lemmatise        & 0.78 & 0.78 \\
\hline
\end{tabular}
\label{tab:roberta_ablation}
\end{table}


The baseline fine tuned RoBERTa model outperformed Bi-LSTM considerable and Naive Bayes significantly without any engineered features and preprocessing. This further highlights how a stronger model architecture (Transformers vs LSTM) and better pre training model (BERT vs GloVE) can make
engineered features and preprocessing less effective, making it difficult for a simple model to bridge the complexity gap with meticulous feature engineering and preprocessing. 

\section{Discussion}

\subsection{Result Analysis: Naive Bayes Hyperparameter Tuning}

\textbf{N-gram Range:} Based on a random search, models which do not take into account unigrams (i.e. ranges starting with 1) performed significantly worse than models that do. In our sample of 50 random searches, the best performing model that did not test unigrams had an F1 of 0.56, while the worst performing model that included unigrams had an F1 of 0.63, showing the importance of using unigrams in the model. Intuitively, this is true because many standalone words like `love' and `hate' offer strong sentiment signals, that cannot be captured as well if only bigrams and longer were considered. In addition, the best performing model suggests that an ngram range of (1,2) is more than enough to detect the sentiment in the text, suggesting that the sentiment is largely made up of single words and short phrases (e.g. not good). A further reason may be the relatively small size of the dataset, which means that there are few examples of longer n-grams, making it difficult for the model to learn reliable estimates for them.

\textbf{Negation Window:} From testing, whether the window used was 3, 4, or 5 words long made no difference to the performance of the model. This suggests the lack of long-range negation dependencies in the text. As such, the final model used a 3-word negation window to simplify computation.

\textbf{Alpha:} Regardless of the other parameters, an alpha of 1.0 always yielded the best F1 score, suggesting Laplace Smoothing is the best configuration. This may again be due to the relatively small size of the dataset, which means that there are many rare words that would benefit from smoothing to prevent zero probabilities.

\textbf{Max features:} From testing, 4000 features was the ideal configuration, with a noticeable drop in F1 score in either direction. This suggests that the text corpus is not too complicated, with there being about 4000 frequently occurring words.

\subsection{Result Analysis: Naive Bayes vs. RoBERTa}

By analyzing the log-likelihood contributions of individual tokens in instances where Naive Bayes (NB) misclassified but RoBERTa succeeded, we identified four primary linguistic patterns that challenge the bag-of-words (BoW) assumption.

\begin{enumerate}
    \item \textbf{Spurious Correlations of Neutral Tokens:} 
    Tokens without inherent sentiment, such as \texttt{`day'} ($1.2487$) and \texttt{`downloaded'} ($1.5053$), frequently contribute to positive predictions simply due to their frequent co-occurrence with positive sentiment-bearing words in the training set. While NB captures these statistical correlations, it lacks the context-awareness to recognize when a neutral word is used in a negative phrase. For example, in the phrase ``worst day ever,'' the token \texttt{'day'} should not contribute to a positive classification. This highlights the fundamental limitation of the BoW assumption, which fails to capture the compositional meaning of words.

    \item \textbf{Sentiment Inversion and Sequential Logic:} 
    NB cannot account for shifts in sentiment polarity within a single document. For instance, the tweet ``was so excited to eat the watermelon... and it was terrible'' was misclassified as positive. This was driven by strong contributions from \texttt{`so excited'} ($2.6367$) and \texttt{`excited'} ($2.2254$), which outweighed the negative contribution of \texttt{`terrible'} ($-3.1951$). 
    
    In contrast, RoBERTa correctly identified the negative sentiment. Its architecture allows it to model the sequential structure of the sentence, recognizing that the initial excitement serves only to intensify the subsequent disappointment. To a transformer-based model, the later tokens effectively "gate" or modify the importance of the preceding ones—a capability NB lacks.

    \item \textbf{Polysemy and Contextual Disambiguation:} 
    NB struggles with polysemous words (words with multiple meanings). For example, the token \texttt{`stuffed'} ($-1.7135$) was associated with negative sentiment in the training data (e.g., ``stuffed nose'' or ``stuffed from overeating''). However, in the tweet ``why would they take a photo with stuffed animals?! that's pretty funny,'' \texttt{`stuffed'} is used neutrally to describe a toy. 
    
    RoBERTa`s self-attention mechanism allows it to disambiguate \texttt{`stuffed'} by attending to neighboring tokens like \texttt{`animals'} and \texttt{`photo'}. This confirms that while NB treats each token as an isolated feature, RoBERTa constructs a dynamic representation based on the specific semantic environment.

    \item \textbf{Sensitivity to Sentence Length and Sparsity:} 
    As a generative model that aggregates token-level probabilities, NB struggles with short sentences or those with sparse in-vocabulary features. In the case of the tweet ``Hahaha! Alright,'' the model predicted a neutral label despite the presence of positive tokens like \texttt{`hahaha'} ($1.0139$) and \texttt{`alright'} ($0.9255$). Because the cumulative sum of these contributions failed to cross the classification threshold, the model defaulted to neutral. This demonstrates that NB requires a higher "density" of evidence compared to RoBERTa, which can derive strong confidence from minimal but high-quality contextual cues.
\end{enumerate}
\section{Conclusion}

In this project, we investigated tweet sentiment classification using Naive Bayes, RoBERTa, and LSTM models. Exploratory analysis showed that textual content was far more informative than the available metadata, leading us to focus on tweet-level preprocessing and modelling. Among the saved results, Naive Bayes achieved about 0.69 accuracy and 0.70 macro-F1 after threshold calibration, demonstrating that a simple model can remain competitive with strong preprocessing. Future work should complete the evaluation of the neural models and compare their gains against their additional computational cost.

Our experiments suggest that preprocessing matters substantially for tweet sentiment classification, especially for simpler models. Handling negation, preserving sentiment-bearing tokens, and calibrating thresholds allowed Naive Bayes to reach a reasonably strong macro-F1 despite its simplicity. This shows that careful text normalization and representation can partially compensate for the limitations of bag-of-words models.

At the same time, the architecture choices in the notebook reflect the expectation that neural models should perform better when contextual understanding is required. RoBERTa, in particular, should be more robust to contextual polarity shifts and compositional meaning. LSTM provides a useful middle ground, offering sequential modelling without the full computational cost of a transformer.

One limitation of the present notebook is that the saved outputs are incomplete for the neural models. This makes it difficult to perform a full empirical comparison directly from the notebook snapshot alone. In the final report, the analysis should therefore distinguish clearly between implemented methods and fully verified numerical results.
%%%%%

\clearpage

\bibliography{refs}
\bibliographystyle{acl_natbib}

\section*{Acknowledgements}
We would like to thank our professors Min-Yen Kan and Christian Von Der Weth, as well as our supervisor Esther Gan for their continued support and guidance throughout the project.

This document has been adapted from the ACL Rolling Review Template (ACL ARR) by Min-Yen Kan.  You may find the original template, which NUS has also contributed to in the past, here: \url{https://www.overleaf.com/latex/templates/acl-rolling-review-template/jxbhdzhmcpdm}.

\section*{Statement of Independent Work}

\noindent 1A. Declaration of Original Work. By entering our Student IDs below, we certify that we completed our assignment independently of all others (except where sanctioned during in-class sessions), obeying the class policy outlined in the introductory lecture. In particular, we are allowed to discuss the problems and solutions in this assignment, but have waited at least 30 minutes by doing other activities unrelated to class before attempting to complete or modify our answers as per the class policy. \\
During coding, LLM models such as Gemini 3.1 Pro and GPT-5.3-Codex were used for code generation, debugging, and optimization. This was done with either the chat feature with tools such as Github Copilot, or with the autocomplete feature in code editors such as Visual Studio Code.
LLMs such as Gemini 3 were also used in the process of writing the report with regards to concept clarification and language/formatting improvement. Snippets of the conversation can be found at \url{https://gemini.google.com/share/63c0ce5b60f4} and \url{https://gemini.google.com/share/7abf536f3366}. 
\noindent Signed, 
[A0251997L, e0957585@u.nus.edu],
[A0252459X, e0958047@u.nus.edu],
[A0257425B, e0968397@u.nus.edu],
[A0252829W, e0958417@u.nus.edu],
[A0252183J, e0957771@u.nus.edu]

\appendix

\section{Hyperparameter Tuning Data for Naive Bayes}
The results of the tuning for Naive Bayes can be found at \url{https://drive.google.com/file/d/1iOQ9JErCyhzH8Y9KnnDg0HTNA0mGL559/view?usp=sharing}.

In each file, results can be interpreted as follows:

\begin{itemize}
\item \texttt{model\_type:} Type of model, NB = Naive Bayes
\item \texttt{model\_\_alpha:} Smoothing parameter for the Naive Bayes model (1.0 indicates Laplace smoothing)
\item \texttt{preprocessor\_\_lemmatise:} Boolean indicating if words were reduced to their dictionary base form
\item \texttt{preprocessor\_\_lowercase:} Boolean indicating if all text was converted to lowercase
\item \texttt{preprocessor\_\_negation\_scoping:} Boolean indicating if the model identifies words affected by ``not" or other negations
\item \texttt{preprocessor\_\_negation\_window:} The number of words following a negation term that are marked as negated
\item \texttt{preprocessor\_\_normalize\_expressive
\_lengthening:} Boolean indicating if repeated characters (e.g., ``loooove") were truncated
\item \texttt{preprocessor\_\_remove\_emoticon:} Boolean indicating if emoticons were deleted from the text
\item \texttt{preprocessor\_\_remove\_number:} Boolean indicating if numeric digits were removed
\item \texttt{preprocessor\_\_remove
\_stopword:} Boolean indicating if common filler words (e.g., ``the", ``is") were removed
\item \texttt{preprocessor\_\_remove\_url:} Boolean indicating if web links were stripped from the data
\item \texttt{preprocessor\_\_translate
\_emoticon:} Boolean indicating if emoticons were converted into text descriptions
\item \texttt{vectorizer\_\_max\_features:} The maximum number of top words/phrases kept based on frequency
\item \texttt{vectorizer\_\_min\_df:} The minimum document frequency required for a term to be included
\item \texttt{vectorizer\_\_ngram\_range:} Length of n-grams used in vectorization, as a range of (min\_n, max\_n)
\item \texttt{vectorizer\_\_vectorizer\_type:} The method of feature extraction, either TF-IDF or Count Vectorizer
\item \texttt{mean\_macro\_f1:} The average F1 score across all classes, treating all classes equally regardless of frequency
\end{itemize}
\end{document}